%!TEX root = ../thesis.tex

\chapter{Introduction and Project Transition}
\label{chap:introduction}

\section{Reproducibility question}

High-resolution shock-capturing (HRSC) finite-volume solvers increasingly run on
graphics processing units (GPUs), whose bandwidth suits stencil updates that are
bandwidth- rather than arithmetic-limited \citep{bard_dorelli_2014}; single
precision (binary32, fp32) halves the bytes moved per state against double
(binary64, fp64). Production codes keep revisiting fp32
and mixed precision
\citep{brogi_etal_2024,wang_xia_chen_2025,higham_mary_2022,vanaEtAl2017singlePrecision},
yet the precision is usually chosen once for a whole code and its effect seldom
measured for the scheme and problem in hand. Ideal magnetohydrodynamics (MHD)
sharpens the point: its thin current sheets and near-degenerate wave structure
are places where a rounding-sized perturbation need not stay that size.

Naming a second-order MUSCL--Hancock scheme with an HLL or HLLD Riemann solver
fixes neither the precision, the compiler's floating-point transformations, the
branch taken at zero wave speed, nor the processor type, so two implementations
of one published scheme may differ. What matters is whether that difference
stays below the discretisation or reference scale, or moves a resolved feature
\citep{plesser2018reproducibility,sandveEtAl2013reproducible}: IEEE formats
standardise operations and rounding, yet finite precision and reordered
reductions still change a computed path
\citep{ieee754_2019,goldberg_1991}.

\section{How Report 1 informed Report 2}
\label{sec:ch1-transition}

Report~1 provided the controlled Euler baseline. Across five cases its matched
strict-IEEE HLLC CPU and GPU outputs agreed at $L_1=L_\infty=0$ and
$\mathrm{ULP}_{\max}=0$, and direct fp32--fp64 differences stayed below the
available reference scale, the largest density reference-scaled ratio reaching
$1.30\times10^{-4}$; fast-math compilation and the HLLC--Rusanov choice each
moved the saved state. Those results redirected the project. Once the
arithmetic was matched the processor type produced no measurable saved-state
difference, so what remained worth measuring was which declared settings decide
whether it stays that way. Report~2 therefore keeps the device
paths matched, ports HLL to CUDA, and varies build semantics and the HLL/HLLD
choice deliberately, treating the device \emph{flag} rather than the device as
the variable. Section~\ref{sec:ch2-priorities} gives the full map, including the
planned move to ideal MHD.

\section{Research questions and scope}

The study asks four questions:

\begin{enumerate}
  \item\label{rq:mhd-validation} How far do numerical-reference comparisons and
        physical-property checks validate the CPU ideal-MHD solvers and the
        tested GPU Harten--Lax--van Leer (HLL) path?
  \item\label{rq:controlled-variation} For matched configurations, how do
        precision, build semantics, Riemann solver and branch rule, processor
        type, thread count and Courant--Friedrichs--Lewy (CFL) number alter the
        stored solution, and what performance follows?
  \item\label{rq:scale-time-mca} How does the fp32--fp64 discrepancy vary with
        resolution and observation time, and how does it compare with the
        spread from Monte Carlo arithmetic (MCA), which rounds operations
        randomly at a chosen virtual precision?
  \item\label{rq:reproducibility-record} Which records are needed to reproduce
        the tested claims?
\end{enumerate}

These are answered for the Brio--Wu, Orszag--Tang and Kelvin--Helmholtz
ideal-MHD cases, with Sod and Liska--Wendroff Configuration~3 as a compact
Euler continuity comparison.

\section{Contributions and report structure}

The report contributes an ideal-MHD extension of the Euler framework, with
generalised Lagrange multiplier (GLM) divergence control, CPU HLL and HLLD
paths, a tested CUDA HLL implementation, and a controlled experiment pipeline
recording effective build semantics and run metadata throughout.
Chapters~\ref{chap:development}--\ref{chap:variation-results} give those
additions, the controlled design, the validation limits and the variation
results; Chapters~\ref{chap:reproducibility-discussion}
and~\ref{chap:conclusions} the discussion and conclusions.
